\subsection{Fenwick Tree / BIT + Consultas Offline (distintos en un rango)}
\label{alg:3-1}

\graybold{Preguntas guía:}

\begin{itemize}
\item ¿Necesito sumas o conteos acumulados sobre un rango que cambia con actualizaciones puntuales?
\item ¿Los updates son de un solo índice y las queries son sobre prefijos o rangos [l, r]?
\item ¿Si no hay updates "online", puedo resolver offline ordenando las queries (p. ej. por su extremo derecho)?
\end{itemize}

\graybold{Utilidad:} Mantener sumas de prefijo (o cualquier operación invertible) sobre un arreglo con actualizaciones puntuales, respondiendo consultas de rango en tiempo logarítmico. Es la base de técnicas como conteo de elementos distintos por rango, inversiones en un arreglo y frecuencias acumuladas.

\graybold{Complejidad:} Construcción \bigO{n}. Update puntual \bigO{\log n}. Query de prefijo \bigO{\log n}. Espacio \bigO{n}.

\graybold{Fundamento teórico:} Cada posición i del arreglo interno almacena la suma de un bloque de tamaño (i \& -i), el bit menos significativo de i. Al actualizar se "sube" sumando (i \& -i); al consultar un prefijo se "baja" restando (i \& -i). Esto aprovecha la representación binaria de los índices para cubrir cualquier prefijo con solo \bigO{\log n} bloques.~\cite{fenwick1994}

\graybold{Ejercicio aplicado}

(SPOJ DQUERY — D-query) Dado un arreglo de n números y q consultas (i, j), responder para cada una cuántos valores distintos hay en el subarreglo [i, j]. Técnica: procesar las consultas offline ordenadas por extremo derecho r, recorriendo el arreglo de izquierda a derecha y manteniendo en el BIT un "1" solo en la última aparición de cada valor.

\graybold{I/O:} n ≤ 30 000 y q ≤ 200 000, hasta \aprox{}430 000 tokens de entrada. El código lee línea a línea con readline (patrón 1): en el caso máximo tarda \aprox{}0.7s. Si el límite de tiempo aprieta, la lectura total + tokenización (patrón 2) es la alternativa más rápida.

\graybold{Código solución}

\begin{lstlisting}[language=Python]
import sys
 
def solve():
    data = sys.stdin.buffer.readline
    n = int(data())
    nums = list(map(int, data().split()))
    q = int(data())
 
    queries = []
    for i in range(q):
        l, r = map(int, data().split())
        queries.append((r, l, i))       # se ordena por r (offline)
    queries.sort()
 
    bit = [0] * (n + 1)
 
    def update(i, val):
        while i <= n:
            bit[i] += val
            i += i & (-i)               # ultimo bit encendido
 
    def query(i):
        res = 0
        while i > 0:
            res += bit[i]
            i -= i & (-i)
        return res
 
    last = {}                           # ultima posicion vista de cada valor
    ans = [0] * q
    cur = 0
    for r, l, qi in queries:
        while cur < r:                  # avanza el puntero hasta r
            cur += 1
            x = nums[cur - 1]
            prev = last.get(x)
            if prev:
                update(prev, -1)        # desactiva la aparicion anterior
            update(cur, 1)              # activa la aparicion actual
            last[x] = cur
        ans[qi] = query(r) - query(l - 1)
 
    sys.stdout.write("\n".join(map(str, ans)))
 
solve()
\end{lstlisting}
